\clearpage

\chapter{Configuración de un bloque de lectura analógica en un pin}
Este capítulo detalla la configuración del periférico Conversor Analógico a Digital (ADC) para realizar una lectura de voltaje, un proceso esencial para interactuar con sensores analógicos. Dada la arquitectura de doble núcleo de la \textbf{STM32H745ZI}, se asume que la lógica de la aplicación y la gestión del periférico ADC se centralizan exclusivamente en el núcleo principal, el \textbf{Cortex-M7}. \\


\section{Configuración del CubeMx}


Usando una copia del archivo \texttt{.ioc} que establecimos como plantilla base, procedemos a configurar los periféricos necesarios para la adquisición de la señal analógica. Usáremos el ADC1, por lo que, en el panel de Analog del STM32CubeMX, se activó el periférico ADC1 y dentro de sus canales se escogió el IN2 - Single-ended (entrada no diferencial) para recibir la señal analógica. Después de esto, el pin físico asociado \textbf{PF11} ahora está en modo de alta impedancia, listo para recibir la señal del sensor analógico. Fig. \ref{fig:adc1}. \\
Cuando la tarjeta nos inicializó todos los pines por default activó el periférico USB OTG FS/HS (Full-Speed / High-Speed). Dado que no lo utilizaremos, lo desactivamos, como indica la fig. \ref{fig:usb}. \\

\begin{figure}
	\centering
	\includegraphics[width=\linewidth]{images/tlk124/adc1}
	\caption{Selección del periférico \textbf{ADC1} y del canal de lectura \textbf{IN2 - Single-ended}. Al realizar esta configuración, el pin físico asociado (generalmente PA1) queda automáticamente configurado en modo analógico de alta impedancia, listo para la adquisición de voltaje en el núcleo \textbf{Cortex-M7}.}
	\label{fig:adc1}
\end{figure}


\paragraph{Configuraciones avanzadas.} 
\begin{itemize}
	\item Se modificó la opción de generación de código para la librería del ADC de HAL (Hardware Abstraction Layer) a LL (Low Layer - o Capa de Bajo Nivel). El Embedded Coder funciona con la capa LL debido a que ofrece un control más directo sobre los registros del microcontrolador. 
	\item En la columna Function Call, se desactivaron todas las casillas. 
	\item De manera similar, se desactivaron todas las casillas en la columna Visibility para las funciones de inicialización.
\end{itemize}

\begin{figure}
	\centering
	\includegraphics[width=\linewidth]{images/tlk124/configavanzada}
	\caption{Configuración final en la pestaña \textbf{Project Manager}. Se selecciona la librería \textbf{LL} para el ADC y se desactivan todas las casillas en las columnas \textbf{Function Call} y \textbf{Visibility}. }
	\label{fig:configavanzada}
\end{figure}

\subsection{Finalización y guardado}
Una vez completados todos los pasos de configuración de los periféricos (SYS, TIM7, ADC1) y las opciones del proyecto (\textbf{Project Manager}), es necesario guardar el trabajo. Diríjase al menú principal y seleccione File $\rightarrow$ Save. \\

\clearpage

\section{Creación del Modelo en Simulink}
	Una vez que la configuración de hardware en el \textbf{STM32CubeMX} ha sido guardada, el siguiente paso es abrir \textbf{MATLAB} y \textbf{Simulink} para diseñar el modelo de adquisición analógica. \\
Desde la ventana de comandos de MATLAB, escriba el comando \texttt{simulink} o acceda a la pestaña HOME y haga clic en el icono de \textbf{Simulink}. Luego seleccione la opción \textbf{Crear Modelo en Blanco} (\textit{Create Blank Model}) y guarde el archivo. \\


\begin{figure}[h]
	\centering
	\includegraphics[width=0.7\linewidth]{images/tlk124/Ventanainicio}
	\caption{Herramienta para abrir Simulink desde la interfaz de Matlab.}
	\label{fig:ventanainicio}
\end{figure}
\begin{figure}[h]
	\centering
	\includegraphics[width=0.7\linewidth]{images/tlk124/blankmodel}
	\caption{Creación de un \textbf{Modelo en Blanco} dentro de la interfaz de Simulink. Este archivo (\texttt{.slx}) servirá como el lienzo principal para el diseño del algoritmo de adquisición de datos y será la base para la generación de código C con el \textit{Embedded Coder}.}
	\label{fig:blankmodel}
\end{figure}

\paragraph{Configuración del bloque.}
Para interactuar con el periférico que configuramos en el paso anterior, debemos usar el bloque específico proporcionado por el paquete de soporte del \textbf{Embedded Coder}. Accedemos a la librería de bloques de Simulink (fig. \ref{fig:lib}) y navegamos a la sección de \textbf{Embedded Coder Support Package for STM32 Microcontrollers}. Dentro de esta sección,  seleccionamos el Analog to Digital Converter, fig. \ref{fig:adcc}.  \\


\begin{figure}[h]
	\centering
	\begin{minipage}[t]{0.49\linewidth}
	\centering
	\includegraphics[width=\linewidth]{images/tlk124/lib}
	\caption{Vista general del \textbf{Library Browser} de Simulink. Aquí se encuentran todas las librerías disponibles, incluyendo la de los bloques base de Simulink y los paquetes de soporte instalados, como el de \textbf{Embedded Coder} para STM32.}
	\label{fig:lib}
	\end{minipage}%
	\hfill
	\begin{minipage}[t]{0.49\linewidth}
		\centering
	\includegraphics[width=\linewidth]{images/tlk124/adcc}
	\caption{Selección del bloque \textbf{Digital-to-Analog Converter (DAC)} dentro de la librería \textbf{Embedded Coder Support Package for STM32}. Estos bloques actúan como la interfaz directa entre el algoritmo de Simulink y los periféricos de hardware del microcontrolador.}
	\label{fig:adcc}
	\end{minipage}
\end{figure}

La configuración del bloque ADC en Simulink debe reflejar la configuración de hardware realizada en CubeMX para el núcleo M7, fig. \ref{fig:parametrizacionanalogtodc}. \\
\begin{itemize}
	\item En el campo de \textbf{ADC module}, se selecciona el \textbf{ADC1} porque este fue el periférico configurado y activado exclusivamente en el \textbf{CubeMX}; esto asegura que el código generado por Simulink acceda a la interfaz correcta. 
	\item En \textbf{Conversion Group} se elige regular. Debido a que el ADC divide las conversiones en dos grupos: \textbf{Regular} (para lecturas periódicas, como la nuestra) e \textbf{Injected} (para lecturas de alta prioridad y sincronizadas con eventos, implica un orden diferente). Seleccionamos \textbf{Regular} porque solo contamos con un pin.
	\item En \textbf{Trigger mode}, la opción a tratar es la de Trigger and Read, pues para este pin no utilizaremos el DMA.  El procedimiento consiste en que el bloque de Simulink genera un trigger por software para \textbf{iniciar la conversión} del ADC y luego espera a que la conversión termine antes de leer el resultado.
	\item \textbf{Number of conversions}. Como solo configuramos un canal analógico (IN2 - Single-ended) en CubeMX, solo necesitamos una conversión por ciclo de ejecución del modelo para obtener el valor deseado.
	\item \textbf{Output status}. Si esta opción estuviera activada, el bloque de Simulink generaría una señal de salida adicional para indicar el estado del ADC (p. ej., si la conversión falló). 
	\item \textbf{Sample time}. Este valor define la frecuencia con la que el bloque ADC intentará leer los datos. Se utiliza el mismo valor de \textbf{Fixed-step size} del Solver (1e-3 segundos) para asegurar que la lectura se realice a la frecuencia de \textbf{1 kHz} definida para el modelo.
\end{itemize}




\begin{figure}
	\centering
	\includegraphics[width=0.7\linewidth]{images/tlk124/parametrizacionAnalogToDC}
	\caption{Ventana de \textbf{Block Parameters} del bloque ADC en Simulink. Se configura el \textbf{ADC module} a \texttt{ADC1}, el \textbf{Conversion group} a \texttt{Regular}, el \textbf{Trigger mode} a \texttt{Trigger and read}, y el \textbf{Sample time} a \texttt{1e-3} segundos, asegurando que la lectura analógica se ejecute a $1\text{ kHz}$ de forma continua.}
	\label{fig:parametrizacionanalogtodc}
\end{figure}




Una vez que el bloque ADC (Analog-to-Digital Converter) está en el modelo, el siguiente paso es conectar su salida a un bloque terminal para poder acceder a los datos generados. Dando doble clic sobre el espacio de trabajo de nuestro modelo, aparece un menú contextual en el cual podemos buscar el bloque denominado \textbf{Outport}; lo seleccionamos y unimos al bloque de Analog to Digital Converter. \\

\begin{figure}[h]
	\centering
	\begin{minipage}[t]{0.49\linewidth}
	\centering
	\includegraphics[width=\linewidth]{images/tlk124/contextmenu}
	\caption{Menú contextual para seleccionar bloques de las librerías de manera más rápida. Seleccionamos el bloque Outport. }
	\label{fig:contextmenu}
	\end{minipage}%
	\hfill
	\begin{minipage}[t]{0.49\linewidth}
	\centering
	\includegraphics[width=\linewidth]{images/tlk124/outport}
	\caption{Conexión del bloque \textbf{ADC} a un bloque \textbf{Outport}. El \textbf{Outport} es crucial, ya que al generar código, este bloque se traduce en una \textbf{variable global} que contendrá el valor de la lectura analógica (raw value), haciéndola accesible para otros subsistemas o código externo en el núcleo M7.}
	\label{fig:outport}
	\end{minipage}
\end{figure}




\paragraph{Configuración de los parámetros del modelo}
El siguiente paso es acceder a la ventana de \textbf{Configuration Parameters} para establecer los parámetros de tiempo y ejecución del modelo. Para ello, es necesario presionar Ctrl+E. La pestaña en la que se va a trabajar es la de \textbf{Solver} y se configuran según las instrucciones de la tab. \ref{tab:solver_config}. En la fig. \ref{fig:parametrosmodelo} se observa el resultado. \\

\begin{table}[h]
	\centering
	\caption{Configuración del Solucionador (Solver) en Simulink}
	\label{tab:solver_config}
	\begin{tabular}{|l|l|p{6cm}|}
		\hline
		\textbf{Parámetro} & \textbf{Configuración} & \textbf{Justificación} \\
		\hline
		Stop Time & \texttt{inf} & Asegura que el modelo continúe ejecutándose indefinidamente una vez cargado en el microcontrolador (M7), imitando el bucle \texttt{while(1)} de un programa embebido. \\
		\hline
		Type & Fixed-step & El código embebido requiere un tiempo de ejecución determinista y predecible. El paso fijo garantiza que la lógica del modelo se ejecute siempre con un intervalo de tiempo constante. \\
		\hline
		Solver & Discrete & Se selecciona un solucionador discreto porque el código del microcontrolador no maneja ecuaciones diferenciales continuas, sino actualizaciones discretas basadas en el tiempo del \textbf{Timer} (TIM7). \\
		\hline
		Fixed-step size & $1\text{e-}3$ & Este valor define el \textbf{Sample Time} del modelo. El valor de $1\text{ ms}$ (milisegundo) se selecciona para una alta velocidad de ejecución, haciendo que el modelo se ejecute a una frecuencia de $1\text{ kHz}$. \\
		\hline
	\end{tabular}
\end{table}

\begin{figure}
	\centering
	\includegraphics[width=0.7\linewidth]{images/tlk124/parametrosmodelo}
	\caption{Configuración de la pestaña \textbf{Solver} dentro de \textbf{Configuration Parameters}. Se establece el \textbf{Stop Time} en \texttt{inf} y se selecciona un \textbf{Fixed-step} discreto de $1\text{e-}3$ segundos. Esta configuración es vital para generar código determinista, asegurando que la lógica del modelo se ejecute a $1\text{ kHz}$ en el microcontrolador.}
	\label{fig:parametrosmodelo}
\end{figure}


Esta ventana es el puente de conexión que le indica a Simulink y al Embedded Coder qué microcontrolador está utilizando y cómo debe generar el código C compatible (fig. \ref{fig:hardwareimplementation}). Para configurarlo, los pasos a seguir son:
\begin{enumerate}
	\item Se selecciona la opción STM32H7xx Based (Dual-core). Se elige esta opción porque la placa Nucleo-H745ZI utiliza el microcontrolador STM32H745ZI, que pertenece a la serie H7 de STM32 y, crucialmente, presenta una arquitectura de doble núcleo (aunque únicamente usemos un npucleo, debe coincidir con la tarjeta que se seleccionó en CubeMX).
	\item Una vez seleccionado el \textbf{Hardware Board} STM32H7xx (Dual-core), Simulink automáticamente se enfoca en el núcleo más potente, el ARM Cortex-M7, para generar el código de la aplicación (el modelo). Dado que en el \textbf{STM32CubeIDE} se configuró el proyecto para que el M7 tuviera el control exclusivo de los periféricos y la aplicación principal, el Embedded Coder sigue esta lógica y genera el código para que se ejecute en el M7. (Además de que en esta ventana no hemos habilitado el uso del otro nucleo).
	\item En el Code Generation System Target File, el campo muestra \verb*|ert.tlc|. Es el archivo de lenguaje de destino (Target Language Compiler) que le indica al Embedded Coder que debe generar código C Este archivo se selecciona en los parámetros de un modelo para indicar que el código generado se optimiza para la velocidad y el uso mínimo de memoria, siendo adecuado para la implementación en hardware embebido. 
\end{enumerate}

\begin{figure}
	\centering
	\includegraphics[width=0.7\linewidth]{images/tlk124/hardwareimplementation}
	\caption{Pestaña \textbf{Hardware Implementation} en \textbf{Configuration Parameters}. Se selecciona \texttt{STM32H7xx Based (Dual-core)} como la tarjeta de hardware, lo que automáticamente dirige la generación de código al núcleo principal \textbf{ARM Cortex-M7} de la placa, necesario para la importación de la configuración del \texttt{.ioc} de CubeMx.}
	\label{fig:hardwareimplementation}
\end{figure}


\begin{figure}
	\centering
	\includegraphics[width=0.6\linewidth]{images/tlk124/initializeboardparameters}
	\caption{Tras la selección de una tarjeta dentro de Hardware Implementation, sale un cuadro de diálogo que induce una espera. Tarda en cambiar el modelo de la tarjeta y es necesario esperar.}
	\label{fig:initializeboardparameters}
\end{figure}

El paso final en la configuración de Simulink es vincular el archivo de configuración de hardware generado por STM32CubeIDE (\texttt{.ioc}) con el modelo de Simulink. Esto permite al \textbf{Embedded Coder} conocer todos los ajustes de periféricos (ADC1, TIM7, GPIO) que se realizaron en el núcleo M7. Al seleccionar \texttt{Browse} nos aparece el explorador de archivos, donde debemos buscar el .ioc.
\begin{figure}
	\centering
	\includegraphics[width=0.7\linewidth]{images/tlk124/seleccionioc}
	\caption{Pestaña \textbf{Hardware Implementation} mostrando la vinculación del archivo de configuración \texttt{.ioc} de CubeIDE. Al cargar este archivo, el \textbf{Embedded Coder} absorbe todos los ajustes de periféricos (ADC1, TIM7) y verifica que la generación de código se realizará para el núcleo \textbf{Cortex-M7}.}
	\label{fig:seleccionioc}
\end{figure}

En el campo de Code Generation, dentro del panel izquierdo de la ventana, seleccionamos en Build process, \textbf{Generate code only}, al igual que nos aseguramos que en Target selection diga: \verb|ert.tlc|, C, y C99 ISO, fig. \ref{fig:codegen}. Al evitar el build automático, nos aseguramos de que Simulink solo invierta tiempo en generar y actualizar los archivos de código C específicos de nuestro modelo, por lo que podría ahorrar tiempo. Mientras no utilicemos el external mode, no tenemos porque construir el archivo. \\

\begin{figure}
	\centering
	\includegraphics[width=0.7\linewidth]{images/tlk124/codegen}
	\caption{Pestaña \textbf{Code Generation}. Se verifica el uso del \textit{System Target File} \texttt{ert.tlc} para código optimizado. Es importante modificar las \textbf{Build options} para \textbf{evitar la compilación automática} (`Build, load and run`), ahorrando tiempo al solo generar los archivos de código del modelo de Simulink.}
	\label{fig:codegen}
\end{figure}

\paragraph{Configuración del Embedded Coder mediante Quick Start}
Para indicar donde queremos que se cree el código de nuestra librería, es necesario cambiar la ruta de generación del código. Para ello, en Matlab, seleccione el icono \texttt{Browse for folder}. Ahora, la ruta de su carpeta debería actualizarse dentro de la barra de búsqueda superior. \\

\begin{figure}
	\centering
	\includegraphics[width=0.7\linewidth]{images/tlk124/broiwseforfolder}
	\caption{Icono del Browse for folder. Para cambiar la ruta donde se guarda la librería.}
	\label{fig:broiwseforfolder}
\end{figure}

\begin{figure}
	\centering
	\includegraphics[width=0.7\linewidth]{images/tlk124/path}
	\caption{Actualización del path en Matlab.}
	\label{fig:path}
\end{figure}


En la barra de herramientas, se selecciona \textbf{APPS} y busca el Embedded Coder. Se creará una pestaña denominada C CODE donde podrá seleccionar la herramienta de Quick Start, fig. \ref{fig:quickstart}. Las fig. \ref{fig:systemquickstart2}, \ref{fig:outputquickstart12}, \ref{fig:system1quickstart2} y \ref{fig:opqz2} resumen el proceso de las opciones elegidas. \\



\begin{figure}[h]
	\centering
	\begin{minipage}[t]{0.49\linewidth}
	\centering
	\includegraphics[width=\linewidth]{images/tlk124/embeddd}
	\caption{Se selecciona el Embedded Coder dentro de la categoría de apps en nuestra barra de herramientas principal.}
	\label{fig:embeddd}
	\end{minipage}%
	\hfill
	\begin{minipage}[t]{0.49\linewidth}
		\centering
	\includegraphics[width=\linewidth]{images/tlk124/quickstart}
	\caption{Herramienta de Quick Start, necesaria para la generación de código de manera más fácil.}
	\label{fig:quickstart}
	\end{minipage}
\end{figure}



\begin{figure}[h]
	\centering
	\begin{minipage}[t]{0.49\linewidth}
		\centering
		\includegraphics[width=\linewidth]{images/tlk124/system}
		\caption{Pestaña System del QuickStart.}
		\label{fig:systemquickstart2}
	\end{minipage}%
	\hfill
	\begin{minipage}[t]{0.49\linewidth}
		\centering
		\includegraphics[width=\linewidth]{images/tlk124/output}
		\caption{Pestaña Output del Quick Start.}
		\label{fig:outputquickstart12}
	\end{minipage}
\end{figure}



\begin{figure}[h]
	\centering
	\begin{minipage}[t]{0.49\linewidth}
		\centering
		\includegraphics[width=\linewidth]{images/tlk124/deploy}
		\caption{Pestaña Deployment del Quick Start.}
		\label{fig:system1quickstart2}
	\end{minipage}%
	\hfill
	\begin{minipage}[t]{0.49\linewidth}
	\centering
	\includegraphics[width=\linewidth]{images/tlk124/optimiza}
	\caption{Pestaña optimización del Quick Start.}
	\label{fig:opqz2}
	\end{minipage}
\end{figure}


\begin{figure}[h]
	\centering
	\begin{minipage}[t]{0.49\linewidth}
	\centering
	\includegraphics[width=\linewidth]{images/tlk124/generationCode}
	\caption{Ventana que indica que la generación de código fue completada correctamente.}
	\label{fig:generationcode}
	\end{minipage}%
	\hfill
	\begin{minipage}[t]{0.49\linewidth}
		\centering
		\includegraphics[width=\linewidth]{images/tlk124/panelcode}
		\caption{Después de la correcta generación de código, se abre el panel derecho donde se muestran los archivos de la librería generada.}
		\label{fig:panelcode}
	\end{minipage}
\end{figure}



\section{Análisis del código generado}
La salida de la consola dio el siguiente mensaje. \\
\begin{minipage}{\textwidth}
	\begin{lstlisting}[caption={Salida de la consola}, label={lst:iconmod}]
	### Searching for referenced models in model 'MB_ADC_1pin_Sk'.
	### Total of 1 models to build.
	### Starting build procedure for: MB_ADC_1pin_Sk 
	
	### Generating code and artifacts to 'Model specific' folder structure
	### Generating code into build folder: E:\git\WaveLink32\Src\TLK1\ManualBuilding

	{...}
	
	.............................
	### Writing header file MB_ADC_1pin_Sk.h
	### Writing source file MB_ADC_1pin_Sk.c
	### Writing header file rtwtypes.h
	### Writing header file rtmodel.h
	.
	### Writing source file ert_main.c
	### TLC code generation complete (took 42.429s).
	### Saving binary information cache.
	### Generating code from STM32CubeMX project: E:\git\WaveLink32\Src\TLK1\ManualBuilding\MB_ADC_1pin\MB_ADC_1pin.ioc
	Scatterfile written to E:\git\WaveLink32\Src\TLK1\ManualBuilding\MB_ADC_1pin\MDK-ARM\CM7\MB_ADC_1pin.sct
	### Define all interrupt handler as weakly defined in stm32h7xx_it.c
	### Define all interrupt handler as weakly defined in stm32h7xx_it.c 
	
	### Successful completion of code generation for: MB_ADC_1pin_Sk
		
	\end{lstlisting}
\end{minipage}		

En la fig. \ref{fig:libreriagenerada} se observa la librería generada. \\
\begin{figure}[h]
	\centering
	\includegraphics[width=0.7\linewidth]{images/tlk124/libreriagenerada}
	\caption{Librería generada tras usar el Embedded Coder. Nótese que la carpeta se denomina (casi) igual que el nombre del proyecto.}
	\label{fig:libreriagenerada}
\end{figure}

Estos archivos contienen el código de aplicación en tiempo real (\textbf{ERT}) generado por el modelo de Simulink ($\texttt{MB\_ADC\_1pin\_Sk}$). El código está optimizado para la eficiencia en la ejecución y en el uso de memoria RAM, y está diseñado para ejecutarse en el núcleo \textbf{ARM Cortex-M7} del STM32H745ZI. La función principal de este código es inicializar el periférico ADC1 y ejecutar repetidamente la lectura de una muestra analógica a la frecuencia de *sampleo* definida en el modelo.

\vspace{0.3cm}
\noindent \textbf{Descripción de los Archivos Generados:}
\begin{itemize}
	\item \texttt{MB\_ADC\_1pin\_Sk.h}: Es el archivo de cabecera principal. Contiene las declaraciones de las funciones del modelo (\texttt{MB\_ADC\_1pin\_Sk\_initialize} y \texttt{MB\_ADC\_1pin\_Sk\_step}), los prototipos de las estructuras de datos globales (\texttt{DW} y \texttt{ExtY}), y define las inclusiones necesarias como \texttt{stm\_adc\_ll.h} y \texttt{MW\_target\_hardware\_resources.h}.
	\item \texttt{MB\_ADC\_1pin\_Sk.c}: Archivo de código fuente principal. Contiene la implementación de la lógica del modelo y las funciones de inicialización. Aquí se encuentran la función de ejecución periódica (\texttt{MB\_ADC\_1pin\_Sk\_step}) y la función de configuración de *hardware* de bajo nivel (\texttt{SystemCore\_setup}).
	\item \texttt{rtwtypes.h}: Archivo de cabecera estándar de Simulink Coder que define los tipos de datos primitivos específicos del código generado (por ejemplo, \texttt{real\_T}, \texttt{uint32\_T}) para garantizar la portabilidad y compatibilidad del código.
	\item \texttt{ert\_main.c} (Template): Archivo de código fuente de plantilla que, si se genera, contiene el bucle principal (\texttt{main}) y llama a las funciones generadas (\texttt{\_initialize} y \texttt{\_step}) en un sistema \textbf{Embedded Real-Time} (ERT).
\end{itemize}


\paragraph{Descripción de las funciones} 
Función de inicialización del modelo. Es el punto de entrada para el código C y se llama una sola vez al inicio del programa. Se encarga de inicializar la estructura de datos del modelo (rtDW.obj) y de llamar a la función local \texttt{SystemCore\_setup}  para configurar el hardware. Lst. \ref{lst:inicia1}. \\
\begin{minipage}{\textwidth}
	\begin{lstlisting}[caption={Función de inicialización}, label={lst:inicia1}]
	void MB_ADC_1pin_Sk_initialize(void);
	\end{lstlisting}
\end{minipage}		

Función de ejecución periódica. Es el corazón del modelo en tiempo real y debe ser llamada por el scheduler (interrupción del TIM7) cada $1\text{ ms}$. Su tarea es ejecutar la lógica del modelo; en este caso, llama a la función regularReadADCNormal para realizar la conversión y actualizar el valor de salida (rtY.Out1). Lst. \ref{lst:inicia11}. \\

\begin{minipage}{\textwidth}
	\begin{lstlisting}[caption={Función de ejecución periódica}, label={lst:inicia11}]
	void MB_ADC_1pin_Sk_step(void);
	\end{lstlisting}
\end{minipage}		

Función estática (solo accesible localmente). Contiene la lógica de inicialización de bajo nivel del ADC1. Configura el manejador del ADC \texttt{ADC\_Handle\_Init} para modo normal, una sola conversión, trigger y lectura, calibra el ADC, lo habilita y comienza la conversión. Lst. \ref{lst:inicia111}.\\

\begin{minipage}{\textwidth}
	\begin{lstlisting}[caption={Inicialización del ADC1}, label={lst:inicia111}]
	static void SystemCore_setup(stm32cube_blocks_AnalogInputFor *obj);
	\end{lstlisting}
\end{minipage}	

\paragraph{Estructuras}
Estructura de Salidas Externas. Contiene las variables globales correspondientes a los bloques Outport del modelo. En este caso, contiene la variable Out1 (tipo \texttt{uint32\_T} ), que almacena el valor de la lectura RAW del ADC. Lst. \ref{lst:inicia1111}. \\
\begin{minipage}{\textwidth}
	\begin{lstlisting}[caption={Estructura de salidas externas}, label={lst:inicia1111}]
	ExtY rtY;
	\end{lstlisting}
\end{minipage}	
Estructura de Datos de Trabajo. Contiene las señales internas y estados del modelo. Almacena la estructura del objeto (obj) del bloque ADC de Simulink, que es esencial para mantener el estado del periférico. Lst. \ref{lst:inicia11111}. \\
\begin{minipage}{\textwidth}
	\begin{lstlisting}[caption={Estados internos del modelo}, label={lst:inicia11111}]
	DW rtDW;
	\end{lstlisting}
\end{minipage}	


\section{Inclusión del código en STM32CubeIDE}
Se recomienda encarecidamente crear una copia del archivo de configuración \texttt{.ioc} que se generó para la configuración del Embedded Coder. De esta forma, nos permite tener una plantilla de configuración de hardware base. Si es necesario crear un nuevo proyecto similar en el futuro, se puede partir de esta copia, ahorrando tiempo. \\
Al trabajar con Embedded Coder, tuvimos que hacer configuraciones muy específicas (como seleccionar la librería LL para el ADC y desactivar las llamadas a funciones en el \textbf{Project Manager}). Para el proyecto final, donde compilaremos todo el código, es posible que deseemos revertir algunas de estas configuraciones a los valores por defecto que son más comunes en un entorno de desarrollo integrado (como volver a usar la generación del main o llamadas a funciones). La copia nos permite hacer esto sin afectar la configuración original que requiere Simulink. \\
\paragraph{Creación del Proyecto Final a partir de la Configuración Base}
Ahora utilizaremos la copia del archivo \texttt{.ioc} para crear el proyecto final en el IDE:
\begin{itemize}
	\item Abrir STM32CubeIDE: Inicie la aplicación.
	\item Iniciar Nuevo Proyecto: Vaya a \textbf{File} $\rightarrow$ \textbf{New} $\rightarrow$ \textbf{STM32 Project}.
	\item Seleccionar Proyecto desde \texttt{.ioc}: En la ventana del Project Selector, elija la opción \textbf{From Existing .ioc File} y seleccione la copia de su archivo \texttt{.ioc} (p. ej., \texttt{Final\_Project.ioc}). Fig. \ref{fig:configioc}.
\end{itemize}.

\begin{figure}[h]
	\centering
	\includegraphics[width=0.6\linewidth]{images/tlk124/configIOC}
	\caption{Creación de un nuevo proyecto en \textbf{STM32CubeIDE} utilizando la opción \texttt{From Existing .ioc File}. Se selecciona la copia del archivo de configuración base (p. ej., \texttt{MB\_ADC\_1pin.ioc}) para generar la estructura del proyecto final, manteniendo las configuraciones de hardware previas.}
	\label{fig:configioc}
\end{figure}



\paragraph{Ajuste del Temporizador TIM2 y Verificación de la Frecuencia del Reloj}

Para que el código generado por Simulink (con un \textit{sample time} de $\mathbf{1\text{e-}3\text{ s}}$) funcione correctamente, debemos configurar un temporizador adicional (\textbf{TIM2}) que se utilizará para el \textit{timing} interno de la aplicación. \\

Antes de configurar el \textbf{TIM2} con los valores de Prescaler y Periodo, es esencial conocer la frecuencia de reloj a la que opera su bus. Por ello, accedemos a la pestaña \textbf{Clock Configuration} y verificamos que la frecuencia de reloj del bus \textbf{APB1} es de $\mathbf{75\text{ MHz}}$. Fig. \ref{fig:relojtim2}. \\

Con esta frecuencia ($75\text{ MHz}$), podremos calcular los valores exactos de Prescaler y Periodo del TIM2 para generar una interrupción periódica cada $1\text{ ms}$, sincronizándolo con la tasa de muestreo del modelo de Simulink. Igualmente, es necesario activar su interrupción.  \\

\begin{figure}
	\centering
	\includegraphics[width=0.7\linewidth]{images/tlk124/relojTim2}
	\caption{Frecuencia con la que funciona el TIM2. Está defenida como 75MHz.}
	\label{fig:relojtim2}
\end{figure}

\begin{figure}
	\centering
	\includegraphics[width=0.7\linewidth]{images/tlk124/tim3confififii}
	\caption{El valor de $\mathbf{75-1 \text{ MHz}}$ es crucial para calcular los parámetros de Prescaler y Periodo del \textbf{TIM2}, asegurando que el temporizador genere un \textit{tick} cada $1\text{ ms}$ para sincronizarse con la tasa de muestreo de $1\text{ kHz}$ de Simulink.}
	\label{fig:tim3confififii}
\end{figure}

\begin{figure}
	\centering
	\includegraphics[width=0.7\linewidth]{images/tlk124/timinte2}
	\caption{Activación de la interrupción del TIM2.}
	\label{fig:timinte2}
\end{figure}

Al crear el proyecto desde el archivo \texttt{.ioc}, es fundamental asegurarse de que el \textbf{STM32CubeIDE} genere el archivo principal de código, main.c, ya que este es el punto de entrada de la aplicación y donde se realizará la llamada a las funciones del código generado por Simulink. En la vista de configuración del microcontrolador, navegue a la pestaña Project Manager y vuelva a habilitar la casilla de \textbf{Generate Main}. \\

\begin{figure}
	\centering
	\includegraphics[width=0.7\linewidth]{images/tlk124/mainagailol}
	\caption{Pestaña \textbf{Project Manager} del STM32CubeIDE. Es obligatorio verificar que la opción \textbf{Generate main} esté habilitada. Esto asegura que el IDE cree el archivo \texttt{main.c}, esencial para albergar el bucle infinito \texttt{while(1)} y las llamadas a las funciones de inicialización y paso de Simulink.}
	\label{fig:mainagailol}
\end{figure}

Una vez creado el proyecto final a partir de la copia del archivo .ioc, es necesario revertir una configuración que se había establecido específicamente para el Embedded Coder de Simulink: la desactivación de la visibilidad de las funciones. Navegue a la pestaña Project Manager y active todas las funciones de inicialización relevantes. En el lst \ref{lst:lib99}, se muestra el resultado del main.c tras habilitar esta opción. \\
\begin{figure}
	\centering
	\includegraphics[width=0.7\linewidth]{images/tlk124/visibilidad}
	\caption{Reversión de la configuración en \textbf{Project Manager}. Se activa la columna \textbf{Visibility} para todas las funciones de inicialización. Esto asegura que los prototipos de las funciones sean accesibles y puedan ser llamadas desde el archivo \texttt{main.c}.}
	\label{fig:visibilidad}
\end{figure}

\begin{minipage}{\textwidth}
	\begin{lstlisting}[caption={En el main, ya aparecen los prototipos e inicialización de los periféricos }, label={lst:lib99}]
		
	void SystemClock_Config(void);
	static void MX_GPIO_Init(void);
	static void MX_USART3_UART_Init(void);
	static void MX_ADC1_Init(void);
	static void MX_TIM2_Init(void);
		
	{...} Dentro del int main(void)
	MX_GPIO_Init();
	MX_ETH_Init();
	MX_USART3_UART_Init();
	MX_ADC1_Init();
	MX_TIM2_Init();
	
	\end{lstlisting}
\end{minipage}	



En la carpeta de Drivers del proyecto general, copiamos la carpeta generada por el EmbeddedCoder. Eliminamos todos los archivos que no terminen con e\verb|.c| y \verb|.h|. Fig. \ref{fig:sourceheader}. \\
\begin{figure}
	\centering
	\includegraphics[width=0.7\linewidth]{images/tlk124/sourceheader}
	\caption{Se mantuvieron solamente los archivos que terminan con .c y.h de la librería generada por EmbeddedCoder.}
	\label{fig:sourceheader}
\end{figure}

Para añadir la inclusión de la carpeta al proyecto del Cortex M7, nos aseguramos de estar dentro del proyecto, seleccionando cualquier archivo perteneciente a este. En la barra de herramientas superior, donde esta la ficha de \textbf{Project}, desplegamos el menú y seleccionamos Properties. \\
En ese cuadro de diálogo, accedemos a la categoría de MCU/MPU GCC Compiler, seleccionamos Include Paths y añadimos la ruta relativa a la carpeta. Aplicamos cambios y cerramos la ventana. \\
\begin{figure}[h]
	\centering
	\begin{minipage}[t]{0.49\linewidth}
	\centering
	\includegraphics[width=\linewidth]{images/tlk124/rutacarpeta}
	\caption{Ruta de la carpeta añadida al path.}
	\label{fig:rutacarpeta}
	\end{minipage}%
	\hfill
	\begin{minipage}[t]{0.49\linewidth}
		\centering
	\includegraphics[width=\linewidth]{images/tlk124/includepathcompiler}
	\caption{Inclusión de la ruta al proyecto del CM7}
	\label{fig:includepathcompiler}
	\end{minipage}
\end{figure}

Para observar los cambios realizados, damos clic derecho en el nombre del proyecto y seleccionamos la opción de refresh. Fig. \ref{fig:refresh}. \\

\begin{figure}
	\centering
	\includegraphics[width=0.5\linewidth]{images/tlk124/refresh}
	\caption{Opción de refresh para actualizar el index del archivo. Se selecciona el nombre del proyecto.}
	\label{fig:refresh}
\end{figure}

Además de añadir los archivos generados por la librería, es necesario juntar los archivos con extensión \verb|.c| dentro de la carpeta \verb|Src|. Para una mejor organización, creamos el folder \texttt{Src\_AdcLib}. \\

\begin{figure}[h]
	\centering
	\begin{minipage}[t]{0.49\linewidth}
	\centering
	\includegraphics[width=\linewidth]{images/tlk124/headersacomodados}
	\caption{Organización de los archivos header de la librería producida por el embedded coder.}
	\label{fig:headersacomodados}
	\end{minipage}%
	\hfill
	\begin{minipage}[t]{0.49\linewidth}
	\centering
	\includegraphics[width=\linewidth]{images/tlk124/source}
	\caption{Organización de los source files de la librería producida por el embedded coder.}
	\label{fig:source}
	\end{minipage}
\end{figure}

Dentro de \verb|Src_AdcLib| buscamos la plantilla de main que brinda el generador de código. Se denomina \verb|ert_main.c|. Para evitar conflictos con variables, es necesario desactivar  el archivo dentro de la construcción. Al archivo, le damos clic derecho, seleccionamos properties y dentro de la categoría de \texttt{C/C++ Build} $\rightarrow$ \texttt{Settings} $\rightarrow$ \texttt{Exclude resource from build}. Fig. \ref{fig:exclusion}. Al guardar los cambios, ahora aparece inhabilitado el archivo, con su icono cruzado por una diagonal. \\


\begin{figure}
	\centering
	\includegraphics[width=0.7\linewidth]{images/tlk124/Exclusion}
	\caption{Exclusión del archivo ert\_main.c}
	\label{fig:exclusion}
\end{figure}

\begin{figure}
	\centering
	\includegraphics[width=0.7\linewidth]{images/tlk124/deshabilitado}
	\caption{Deshabilitación del archivo del ert\_main.c}
	\label{fig:deshabilitado}
\end{figure}


\clearpage

\section{El Archivo Guía \texttt{ert\_main.c}}

El archivo \texttt{ert\_main.c} es una plantilla de código fuente de ejemplo, vital para comprender cómo debe ser el punto de entrada de la aplicación en nuestro \texttt{main.c} final. Nos proporciona una guía clara sobre \textbf{Librerías y Funciones a Incluir:} Indica qué archivos de cabecera (\texttt{.h}) generados por Simulink (\texttt{MB\_ADC\_1pin\_Sk.h}) y qué librerías del sistema deben incluirse en nuestro \texttt{main.c}. \\
Por lo tanto, la estructura del \texttt{ert\_main.c} servirá de \textbf{plantilla y guía} para las modificaciones que realizaremos en el archivo \texttt{main.c} generado por STM32CubeIDE. \\

\begin{minipage}{\textwidth}
	\begin{lstlisting}[caption={Inclusión de las librerías en el ert\_main.c}, label={lst:lib1}]
		#include "MB_ADC_1pin_Sk.h"
		#include "rtwtypes.h"
		#include "MW_target_hardware_resources.h"
	\end{lstlisting}
\end{minipage}	

\begin{minipage}{\textwidth}
	\begin{lstlisting}[caption={Sintentización de la función rt\_OneStep que brinda el EmbeddedCoder}, label={lst:lib2}]
	void rt_OneStep(void)
	{
	MB_ADC_1pin_Sk_step();
	/* Get model outputs here */
	}
	\end{lstlisting}
\end{minipage}	

\begin{minipage}{\textwidth}
	\begin{lstlisting}[caption={En el main, se creó la variable de modelBaseRate que nos recuerda el valor asignado en la configuración que hicimos en Simulink. }, label={lst:lib3}]
		
	int main(int argc, char **argv){
	float modelBaseRate = 0.001;
	...
	}
	\end{lstlisting}
\end{minipage}	

En lst. \ref{lst:lib4} se muestra la inicialización de los periféricos junto con el desarrollado durante este capítulo. Dentro de esto, lo importante es todo lo relacionado al ADC y el momento en que activa las interrupciones del sistema embebido. \\

\begin{minipage}{\textwidth}
	\begin{lstlisting}[caption={Inicialización de los periféricos. Equivalente a la configuración que nos dan por default en el main.c y añade las funciones de inicialización del ADC. }, label={lst:lib4}]

// Start Peripheral initialization imported from STM32CubeMX project;
HAL_Init();
SystemClock_Config();
PeriphCommonClock_Config();
MX_GPIO_Init();
MX_ETH_Init();
MX_USART3_UART_Init();
MX_ADC1_Init();
configureADCChannelPreselectionRegister(ADC1,1 ,0);

// End Peripheral initialization imported from STM32CubeMX project;
((void) 0);
MB_ADC_1pin_Sk_initialize();
__disable_irq();
ARMCM_SysTick_Config(modelBaseRate);
runModel = true;
__enable_irq();
__enable_irq();
	\end{lstlisting}
\end{minipage}	


\section{Corrección del código}
Los segmentos de código mencionados en la sección anterior, derivados del archivo guía \texttt{ert\_main.c}, deben ser añadidos cuidadosamente en las diferentes secciones que permiten al usuario escribir su código dentro del archivo \texttt{main.c} generado por STM32CubeIDE.

El IDE utiliza \textbf{comentarios especiales} para delimitar las áreas donde el usuario puede insertar código personalizado sin que este sea eliminado o sobrescrito durante futuras regeneraciones del proyecto basadas en el archivo \texttt{.ioc}.

Las etiquetas de usuario son:
\begin{itemize}
	\item \texttt{/* USER CODE BEGIN ... */}
	\item \texttt{/* USER CODE END ... */}
\end{itemize}

Es crucial insertar cada segmento de código en su ubicación correcta dentro de estas etiquetas para garantizar la integridad y portabilidad del proyecto.
\vspace{0.2cm}

\begin{minipage}{\textwidth}
	\begin{lstlisting}[caption={Inclusión de la librería dentro del main.c }, label={lst:lib5}]
		/* Private includes ----------------------------------------------------------*/
		/* USER CODE BEGIN Includes */
		#include "MB_ADC_1pin_Sk.h"
		#include "rtwtypes.h"
		#include "MW_target_hardware_resources.h"
		/* USER CODE END Includes */
	\end{lstlisting}
\end{minipage}	

\begin{minipage}{\textwidth}
	\begin{lstlisting}[caption={Inclusión de la función proporcionada dentro del user code begin.}, label={lst:lib6}]
		/* USER CODE BEGIN 0 */
		void rt_OneStep(void)
		{
			MB_ADC_1pin_Sk_step();
			/* Get model outputs here */
		}
		/* USER CODE END 0 */
	\end{lstlisting}
\end{minipage}	

\begin{minipage}{\textwidth}
	\begin{lstlisting}[caption={Llamada a la función de inicialización y la habilitación del TIM2 con su correspondiente interrupción. }, label={lst:lib7}]
	 /* Initialize all configured peripherals */
	MX_GPIO_Init();
	MX_ETH_Init();
	MX_USART3_UART_Init();
	MX_ADC1_Init();
	MX_TIM2_Init();
	/* USER CODE BEGIN 2 */
	
	MB_ADC_1pin_Sk_initialize();
	__disable_irq();
	HAL_TIM_Base_Start_IT(&htim2);
	__enable_irq();
	__enable_irq();
	/* USER CODE END 2 */
	\end{lstlisting}
\end{minipage}	

\begin{minipage}{\textwidth}
	\begin{lstlisting}[caption={Modificación del callback de la interrupción de los timers}, label={lst:lib7}]
	void HAL_TIM_PeriodElapsedCallback(TIM_HandleTypeDef *htim)
	{
		/* USER CODE BEGIN Callback 0 */
		
		/* USER CODE END Callback 0 */
		if (htim->Instance == TIM7) {
			HAL_IncTick();
		}
		/* USER CODE BEGIN Callback 1 */
		if (htim->Instance == TIM2)
		{
			rt_OneStep();
		}
		
		/* USER CODE END Callback 1 */
	}
	\end{lstlisting}
\end{minipage}	

\paragraph{Añadir librerías}. Al intentar compilar el código nos aparecen los siguientes errores en la consola. Esto es, porque son necesarios unos archivos que se encuentran dentro de la carpeta de instalación de Matlab. Dentro de la carpeta \verb|C:ProgramData\MATLAB\SupportPackages\R2025b\toolbox\shared\supportpackages\stm32\src| buscamos el archivo \texttt{stm\_adc\_ll.c} y del \verb|C:\ProgramData\MATLAB\SupportPackages\R2025b\toolbox\shared\supportpackages\stm32\include| copiamos el archivo a nuestra carpeta que hicimos dentro de Drivers. \\

\begin{minipage}{\textwidth}
	\begin{lstlisting}[caption={Errores de consola}, label={lst:lib8}]
	../../Drivers/MB_ADC_1pin_Sk_ert_rtw/MB_ADC_1pin_Sk.h:29:10: fatal error: stm_adc_ll.h: No such file or directory
	29 | #include "stm_adc_ll.h"

	\end{lstlisting}
\end{minipage}	

\begin{figure}
	\centering
	\includegraphics[width=0.7\linewidth]{images/tlk124/src}
	\caption{Se incluyó el archivo stm\_adc\_ll.h dentro de nuestra carpeta de MB\_ADC\_1pin\_Sk\_ert\_rtw}
	\label{fig:src}
\end{figure}

\begin{figure}
	\centering
	\includegraphics[width=0.7\linewidth]{images/tlk124/src2}
	\caption{Dentro de la carpeta Src\_AdcLib se añadió el archivo stm\_adc\_ll.c}
	\label{fig:src2}
\end{figure}

Igualmente, es necesario añadir las librerías que estos nuevos archivos llaman dentro de su código. \\

\begin{minipage}{\textwidth}
	\begin{lstlisting}[caption={Errores de consola}, label={lst:lib54}]
	./../Drivers/MB_ADC_1pin_Sk_ert_rtw/stm_adc_ll.h:8:10: fatal error: mw_stm32_types.h: No such file or directory
	8 | #include "mw_stm32_types.h"
		
	\end{lstlisting}
\end{minipage}	

\begin{figure}
	\centering
	\includegraphics[width=0.7\linewidth]{images/capMacros/libaa}
	\caption{Librería añadida}
	\label{fig:libaa}
\end{figure}


\paragraph{Creación de la macro}. En el archivo de \texttt{stm\_adc\_ll.c} es posible observar que existen varias líneas de código subrayadas de color gris. Esto es debido a que es necesario, la definición de la macro que habilita la funcionalidad de ese bloque de código.\\
Usando las instrucciones descritas en el capítulo de Macros, añadimos esta definició: \verb|MW_ADC1_ENABLED|. \\

\begin{figure}[h]
	\centering
	\includegraphics[width=0.7\linewidth]{images/capMacros/macronecesaria}
	\caption{Se añadió la macro a la configuración del proyecto del chip M7.}
	\label{fig:macronecesaria}
\end{figure}


\clearpage
\subsubsection{Borrado de la rutina de doble núcleo}
Se borraron las líneas de código lst. \ref{borrado}, debido a que únicamente se está utilizando el chip M7 y a la hora de debuggear, se puede quedar atascado esperando una respuesta del núcleo M4 inhabilitado.

\begin{minipage}{\textwidth}
	\begin{lstlisting}[caption={Borrado de la rutina de doble núcleo }, label={lst:borrado}]
		#define DUAL_CORE_BOOT_SYNC_SEQUENCE
		
		#if defined(DUAL_CORE_BOOT_SYNC_SEQUENCE)
		#ifndef HSEM_ID_0
		#define HSEM_ID_0 (0U) /* HW semaphore 0*/
		#endif
		#endif /* DUAL_CORE_BOOT_SYNC_SEQUENCE */
		
		{...}
		#if defined(DUAL_CORE_BOOT_SYNC_SEQUENCE)
		int32_t timeout;
		#endif /* DUAL_CORE_BOOT_SYNC_SEQUENCE */
		
		{...}
		
		#if defined(DUAL_CORE_BOOT_SYNC_SEQUENCE)
		// Wait until CPU2 boots and enters in stop mode or timeout
		timeout = 0xFFFF;
		while((__HAL_RCC_GET_FLAG(RCC_FLAG_D2CKRDY) != RESET) && (timeout-- > 0));
		if ( timeout < 0 )
		{
			Error_Handler();
		}
		#endif //DUAL_CORE_BOOT_SYNC_SEQUENCE
		
		{...}
		
		#if defined(DUAL_CORE_BOOT_SYNC_SEQUENCE)
		When system initialization is finished, Cortex-M7 will release Cortex-M4 by means of
		HSEM notification
		HW semaphore Clock enable
		__HAL_RCC_HSEM_CLK_ENABLE();
		Take HSEM
		HAL_HSEM_FastTake(HSEM_ID_0);
		Release HSEM in order to notify the CPU2(CM4)
		HAL_HSEM_Release(HSEM_ID_0,0);
		wait until CPU2 wakes up from stop mode
		timeout = 0xFFFF;
		while((__HAL_RCC_GET_FLAG(RCC_FLAG_D2CKRDY) == RESET) && (timeout-- > 0)){
			asm("nop");
		}
		if ( timeout < 0 )
		{
			Error_Handler();
			
		}
		#endif  DUAL_CORE_BOOT_SYNC_SEQUENCE */
		
	\end{lstlisting}
\end{minipage}	


\section{Resultado}
\begin{minipage}{\textwidth}
	\begin{lstlisting}[caption={Mensaje de que el código no tiene errores de compilación}, label={lst:lib56}]
		Finished building: default.size.stdout
		
		Finished building: MB_ADC_1pin_CM7.list
		
		03:39:28 Build Finished. 0 errors, 1 warnings. (took 8s.59ms)
	\end{lstlisting}
\end{minipage}	

